% !TeX root = ../../Book5.tex
% 纲要标题：### 10.2 反射函子
% 纲要位置：计划纲要.md，第 3407 行。

\section{反射函子}
\label{b5:sec:1002}

% TODO: 按以下纲要要点撰写本节。

% 纲要内容（仅作写作计划，尚非教材正文）：
%
% * 汇点和源点反射
% * 对不可分解表示的作用
% * 与简单反射的对应
% * 明确汇点反射采用核、源点反射采用余核；给出排除相应简单表示后互为逆的范围，不把反射函子当作整个模范畴上的无条件等价
%

待撰写。
